\scanpage{80}
\mspage{II. 28}
% archive sheet number: 79
% lines 1–4
Alors, si \deleted{\unsure{$\ell\nmid$}} $\ell\nmid N$, \deleted{on aura $\mu_\ell(k^*)$} $\pi_0$ opère non trivialt sur le \deleted{\unsure{groupe}} \inserted{$\mathbb F_\ell$-espace} vectoriel de dim 1 $\mu_\ell(\bar k^*)^{\unsure{\otimes1|\otimes-1}}$, donc sur son dual $V$, donc \deleted{la relation} $V^{\pi_0}=0$, ce qui implique
% line 5
\[H^1(\bar X,\mathbb Z/\ell\mathbb Z)^\pi\subset H^1(\bar X',\mathbb Z/\ell\mathbb Z)^\pi\quad\text{si }\ell\deleted{\unsure{\mid}}\nmid N\tag*{(4.2.3)}\]
% lines 6–8
\deleted{Or on a, si $A$ est la jacobienne de $X'$}
\[H^1(\bar X',\mathbb Z/\ell\mathbb Z)\simeq\deleted{\unsure{\operatorname{Hom}({}_\ell A(\bar k))|T_\ell(A(\bar k))}}{}_{\ell}A(\unsure{\bar k^*|\bar k})^\vee\deleted{\unsure{\otimes\mu_\ell(\bar k)^{\otimes1}}},\]
\deleted{où le $\vee$ désigne le dual à valeurs dans $\mathbb Z/\ell\mathbb Z$.}
% lines 9–13
D'autre part, \deleted{\unsure{pour le} $H^1(\bar X',\mathbb Z/\ell\mathbb Z)^\pi=0$ signifie} comme dans 2.5 on se ramène, pour la relation \deleted{i.e. ${}_\ell A(\bar k)^{\vee\pi}=0$ i.e. $({}_\ell A(\bar k))_\pi=0$,} $H^1(\bar X',\mathbb Z/\ell\mathbb Z)^\pi=0$ pour presque $\unsure{\forall|\text{tout}}\ell$, au cas où $k$ est fini.
% lines 14–18
\deleted{Alors il est bien connu, que} \deleted{\unsure{Soit $f$}} Soit donc \deleted{\unsure{comme}} $f$ le générateur \inserted{de Frobenius} de \deleted{\unsure{car $p($}} $\pi_0\simeq\widehat{\mathbb Z}$, et soit $V={}_\ell A(\bar k)$. Comme $V$ est fini, $\operatorname{Ker}(1-f)_V$ et $\operatorname{Coker}(1-f_V)$ ont même cardinal, donc \deleted{la relation} \deleted{\unsure{est}}
\[({}_\ell A(\bar k))_\pi=0\Longleftrightarrow({}_\ell A(\bar k))^\pi=0\quad\text{i.e. }{}_\ell A(k)=0.\]
% lines 19–23
Or soit $N'$ l'\deleted{nombre d'élé} ordre du groupe \inserted{fini} $A(k)$. On voit que \deleted{\unsure{si $\ell$}}
\[{}_\ell A(k)=0\quad\text{si}\quad\ell\nmid N'.\]
Par suite, on trouve que si $\ell\nmid NN'$, alors
\[H^1(\bar X^{\unsure{\ell|'}},\mathbb Z/\ell\mathbb Z)^\pi=0,\quad\text{qfd.}\]
% lines 24–27; corollary continued in lower margin
\uline{Corollaire} \inserted{4.3} Soit $X$ un schéma \inserted{normal} géométriquement \deleted{intègre} \inserted{conn. de t.f.} sur un corps fini $\mathbb F_{p^{\unsure{\nu|v}}}$. Alors pour tout nb premier $\ell\neq p$, \inserted{l'hom. surj.} $\pi_1(X,\xi)^{\mathrm{ab}}(\ell)\longrightarrow\pi_1(\mathbb F_p)\simeq\widehat{\mathbb Z}$ \deleted{\unsure{$\mathbb Z$ est libre de rang}} a un noyau fini, et pour presque tout $\ell$, c'est un isomorphisme.
